\chapter{Soziale und kulturelle Erfolgsfaktoren der \DO-Transformation}
\label{ch:social-cultural-factors}

Das Verständnis der förderlichen sozialen und kulturellen Faktoren ist für Führungskräfte und Teams entscheidend, um \DO aktiv zu gestalten, anstatt lediglich auf die erhoffte Wirkung technischer Implementierungen zu warten.
Dieses Kapitel destilliert aus der Forschung die zentralen soziokulturellen Erfolgsfaktoren.

\section{Vertrauenskultur und Fehlerfreundlichkeit}

Vertrauen ist das Fundament, auf dem eine effektive \DO-Kultur aufgebaut wird.
Es ermöglicht die notwendige offene Zusammenarbeit zwischen Entwicklungs"= und Betriebsteams, die traditionell durch Misstrauen und getrennte Zielvorgaben geprägt waren~\cites{mas18,akb22}.
Vertrauen manifestiert sich in der Bereitschaft, sich auf die Kompetenzen der anderen zu verlassen und Verantwortung zu teilen.

Eng damit verbunden ist das Konzept einer fehlerfreundlichen Kultur.
In einer solchen Kultur werden Fehler nicht als individuelles Versagen, sondern als gemeinsame Lernchancen betrachtet.
Statt nach Schuldigen zu suchen, wird eine gemeinsame Analyse der Ursachen durchgeführt, um zukünftige Fehler zu vermeiden.
Sie ermutigt Teammitglieder, Risiken einzugehen, Probleme offen anzusprechen und aus Fehlern zu lernen, anstatt sie zu vertuschen~\cite{san18}.
Dieser Ansatz schafft eine Atmosphäre psychologischer Sicherheit \textendash\ eine Grundvoraussetzung für eine produktive Arbeitskultur.
% TODO Beispiel nennen, das zeigt, dass es solche Arbeitskulturen gibt - oder auf Ideal verweisen?

\section{Interdisziplinäre Kommunikation und Wissensaustausch}

\emph{Kommunikation und Kollaboration} sind die am häufigsten genannten und wohl kritischsten Attribute einer funktionierenden \DO-Kultur~\cite{san18}.
Eine fehlende oder ineffektive Kommunikation wird konsistent als eine der größten Hürden bei der Einführung identifiziert~\cite{akb22}.
Es geht dabei um eine kontinuierliche, offene und transparente Kommunikation über alle Teamgrenzen hinweg.
Ein Schlüsselelement zur Überwindung von Silos ist das kontinuierliche Teilen von Wissen.
In traditionellen IT"=Organisationen ist Wissen oft auf einzelne Abteilungen beschränkt.
\DO fördert hingegen aktiv den Wissensaustausch durch Praktiken wie Pair Programming, interne Tech"=Talks und teamübergreifende Projekte.
Ziel ist es, ein gemeinsames Verständnis für den gesamten Software"=Lebenszyklus zu schaffen und Wissensinseln aufzubrechen~\cites{cup16,aza22,kha23}.

\section{Führung und Empowerment}

Die Rolle der Führung verändert sich in einer \DO-Transformation fundamental.
Traditionelle, auf Kontrolle basierende Führungsstile erweisen sich als hinderlich.
Stattdessen sind moderne Unterstützungsstrukturen gefragt.
Solche Führungskräfte inspirieren ihre Teams, fordern den Status quo heraus und schaffen eine Vision, die über die täglichen Aufgaben hinausgeht~\cite{dor24}.

Die Managementforschung identifiziert in diesem Kontext zwei kritische Erfolgsfaktoren: die explizite Unterstützung durch das Top"=Management und die Ermächtigung (\foreignlanguage{american}{Empowerment}) der Teams.
Die Transformation muss von der obersten Führungsebene getragen, kommuniziert und mit den notwendigen Ressourcen ausgestattet werden.
Gleichzeitig müssen Teams die Autonomie erhalten, Entscheidungen über ihre eigenen Arbeitsprozesse und Werkzeuge zu treffen.
% TODO Quelle? Einseitige, vereinfachte Darstellung: Gute Gründe für streng hierarchische Strukturen, z.B. schnelle Entscheidungen treffen zu können.
Mikromanagement und starre Vorgaben untergraben die Eigenverantwortung, die für kontinuierliche Verbesserung notwendig ist.

\paragraph{}
Diese Erfolgsfaktoren bilden die positive Seite der Transformation.
Doch ebenso wichtig ist das Verständnis für die Hürden, die diesem Wandel im Weg stehen.
Das folgende Kapitel beleuchtet daher die Herausforderungen und Widerstände, die der Einführung von \DO entgegenwirken.